\documentclass[10pt,a4paper,twocolumn]{article}

\usepackage[T1]{fontenc}
\usepackage[utf8]{inputenc}
\usepackage[english]{babel}
\usepackage{lmodern}
\usepackage[margin=22mm]{geometry}
\usepackage{microtype}
\usepackage{amsmath}
\usepackage{booktabs}
\usepackage{enumitem}
\usepackage{xcolor}
\usepackage{needspace}
\usepackage{hyperref}

\definecolor{linkblue}{HTML}{173E66}
\hypersetup{
  colorlinks=true,
  linkcolor=linkblue,
  urlcolor=linkblue,
  citecolor=linkblue,
  pdftitle={OpenMP Parallelization of a Feedforward Neural Network},
  pdfsubject={Implementation and experimental analysis}
}
\setlength{\parindent}{0pt}
\setlength{\parskip}{0.35em}
\setlength{\columnsep}{7mm}
\setlist{nosep,topsep=0.4em,leftmargin=*}
\setlength{\emergencystretch}{2em}
\raggedbottom

\title{\textbf{OpenMP Parallelization of a\\Feedforward Neural Network}\\[0.5em]
  \large Implementation and Experimental Analysis}
\author{Parallel Programming Laboratory\\Universitat Aut\`onoma de Barcelona}
\date{}

\begin{document}
\maketitle

% Complete the abstract with measured results and conclusions after the
% implementation and experimental evaluation are finished.
\begin{abstract}
We investigate OpenMP parallelization of the supplied feedforward neural
network, progressing from a corrected sequential baseline to loop-based and
task-based implementations. The report follows our implementation decisions
and their measured outcomes. Evaluation focuses on execution time, speedup,
efficiency, and strong scaling, with brief checks of neural-network behavior.
Findings and conclusions will be added after the experiments.
\end{abstract}

\section{Experimental approach}
\label{sec:approach}

We parallelize the supplied C program and evaluate its performance as the
thread count increases~\cite{assignment}. After examining dependencies and
establishing a corrected sequential baseline, we develop loop-based
parallelism, then investigate a task-based alternative. The assignment
permits either approach; exploring both is our choice. Comparisons preserve
the workload, compiler settings, and timing boundaries. The original source
remains recoverable in Git.

\subsection{Measurement and scaling}

We briefly check successful execution, finite losses, and consistent
recognition results. Evaluation focuses on parallel performance. We also
examine shared data and synchronization to ensure valid execution.

Following the course notation~\cite{performance-metrics}, let $N$ denote
problem size, $p$ the thread count, $T(N,1)$ our measured sequential time,
and $T(N,p)$ the parallel time. Then
\begin{align}
  S(N,p) &= \frac{T(N,1)}{T(N,p)}, \\
  E(N,p) &= \frac{S(N,p)}{p}.
\end{align}
We use our own sequential baseline on the same hardware and workload.
OpenMP runs with one thread provide an additional overhead comparison;
published timings remain a separate reference.

We evaluate \textbf{strong scaling}: keep the dataset, network, batch size,
and epochs fixed and vary $p\in\{1,2,4,6,12\}$. We will plot time, speedup,
and efficiency against thread count. Ideally, $T(N,p)=T(N,1)/p$,
$S(N,p)=p$, and $E(N,p)=1$. These curves show where additional threads stop
helping. \textbf{Weak scaling} grows the workload with the thread count and
is outside this fixed-workload experiment. We will report repeated timings
and variability, and distinguish measured explanations from hypotheses.

\subsection{Course references}

Table~\ref{tab:reference-times} gives indicative course timings for an
optimized OpenMP implementation~\cite{reference-times}. We have access to
Wilma and Aolin, but the reference source code and exact timing boundaries
are unavailable, limiting direct comparisons.

\begin{table}[htbp]
  \centering
  \caption{Course reference execution times in seconds (10 epochs).}
  \label{tab:reference-times}
  \small
  \setlength{\tabcolsep}{4pt}
  \begin{tabular}{@{}lrrr@{}}
    \toprule
    & Wilma & Wilma & Aolin \\
    Threads & \texttt{nodo.q} & \texttt{new-nodo.q} & \texttt{aolin.q} \\
    \midrule
    Sequential & 40.737795 & 12.095956 & 11.235670 \\
    2 & 22.542747 & 6.873421 & 6.767024 \\
    4 & 11.730357 & 4.296542 & 3.658490 \\
    6 & 8.308845 & 3.474935 & 2.502528 \\
    12 & 5.263955 & 2.555904 & 2.380782 \\
    \bottomrule
  \end{tabular}
\end{table}

\section{Understanding the sequential program}
\label{sec:sequential}

\subsection{Execution flow}

In \texttt{main()}, configuration and initialization precede the timer;
training and testing follow, before the timer stops and final cleanup
begins~\cite{starter-code}. The measured interval includes training,
testing, dataset loading, console output, and input-data cleanup. We
preserve this boundary. Testing uses the trained weights, so training and
testing must remain ordered.

The default network has layers $1024\to500\to210\to135\to10$, with
10 training epochs and batches of 50 patterns. The supplied batch loop
processes 38 complete batches per epoch: 1,900 of the 1,934 training patterns,
after shuffling. We preserve this workload, including omission of the last
34 patterns. Testing uses all 946 test patterns.

\subsection{Baseline corrections}

Before benchmarking, we corrected four implementation issues:
\begin{itemize}
  \item \textbf{Image loading:} a bounded row buffer prevents the string
    terminator from overflowing the pixel array.
  \item \textbf{Random-number generation:} unsigned 32-bit arithmetic avoids
    signed overflow while preserving the usual wraparound sequence.
  \item \textbf{Dataset loading:} training targets are initialized only when
    loading training data, preventing unintended writes during testing.
  \item \textbf{Bias handling:} gradients are reset and applied to hidden
    and output biases, excluding the input layer.
\end{itemize}
Both sequential and parallel versions use these corrections. Updating the
output biases can change recognition results.

\section{Sequential baseline}

After a preliminary execution, we collected five sequential runs on
\texttt{clus07.hpc.local}, in Wilma's \texttt{new-nodo.q} partition.
The node reports an Intel Xeon Gold 5416S, two sockets with 16 physical cores
each, and two hardware threads per core (32 physical cores, 64 logical CPUs).
We used GCC 13.2.1 with \texttt{-Ofast}, linking \texttt{-lm}, and kept the
configuration and executable unchanged between runs.

\begin{table}[htbp]
  \centering
  \caption{Sequential application timings on \texttt{clus07}, five runs.
    Standard deviation is the sample standard deviation.}
  \label{tab:sequential-repetitions}
  \small
  \begin{tabular}{@{}lr@{}}
    \toprule
    Statistic & Time (s) \\
    \midrule
    Mean & 12.343398 \\
    Standard deviation & 0.108668 \\
    Minimum & 12.231889 \\
    Maximum & 12.474166 \\
    Course sequential reference & 12.095956 \\
    \bottomrule
  \end{tabular}
\end{table}

The standard deviation is 0.88\% of the mean. All five runs produced the
same printed epoch losses, a final loss of 0.055671, and 916 correct
classifications out of 946. We use the mean application time,
12.343398~s, as $T(N,1)$ for comparisons on this node and configuration;
\texttt{perf}'s process-wide elapsed time is not used. The individual runs
and environment are retained in the
\href{../../results/sequential_baseline.json}{measurement record}.
Our mean exceeds the course reference by 2.05\%; the cause is undetermined.
Up to 12 threads fit on physical cores; placement will be checked.

\section{Loop-based OpenMP}

\subsection{Forward propagation}

We added \texttt{parallel for schedule(static)} to the neuron loop in
\texttt{forward\_prop()}~\cite{openmp-lecture}. Each iteration reads the
preceding layer's activations, weights, and bias, and writes distinct
\texttt{z[j]} and \texttt{actv[j]} elements. Layers remain ordered, with
loop completion ensuring that all activations are available to the next
layer. The inner sum remains sequential, retaining its accumulation order.
Static scheduling suits the equal-length sums. Repeated parallel regions
and synchronization may limit gains in small layers; this requires measurement.

\subsection{Backpropagation and shared sums}

Different neurons in the supplied \texttt{back\_prop()} add to the same
\texttt{dactv[k]}, making direct neuron-loop parallelization unsafe. We
separated the work: parallel neuron iterations compute \texttt{dz}, bias
gradients, and distinct weight-gradient rows; parallel \texttt{k} iterations
then accumulate individual \texttt{dactv[k]} values in the original neuron
order. Loop completion enforces dependencies; input-layer derivatives are
omitted as before. Each element has one writer per phase. Patterns remain
sequential because they share activation buffers and batch gradients.

\subsection{Gradient reset and weight updates}

Next, we will distribute weight-gradient resets across array elements and
weight updates across independent neuron pairs using \texttt{collapse(2)}.
Both use static scheduling and distinct writes, without reductions. The
small bias loops remain sequential to limit parallel-region overhead.
Gradient reset, batch accumulation, and parameter updates remain ordered.
Benchmarking follows completion of these changes and OpenMP build setup.

\Needspace{7\baselineskip}
\begin{thebibliography}{9}
  \small
  \raggedright
  \bibitem{assignment}
    Universitat Aut\`onoma de Barcelona,
    \textit{Parallel Programming Laboratory Project},
    2026--2027, Sections III--IV.
    \href{../../resources/CAP_Pra_ctica_NN_26_27.pdf}{Repository copy}.
  \bibitem{reference-times}
    Universitat Aut\`onoma de Barcelona,
    \textit{FeedForward Neural Networks: Reference Time -- SEQ + OpenMP},
    October 2026.
    \href{../../resources/ReferenceTimesOpenMP26_27.pdf}{Repository copy}.
  \bibitem{performance-metrics}
    Universitat Aut\`onoma de Barcelona,
    \textit{Lecture 3: Introduction to Parallelism},
    slides 20--22, 26--30.
  \bibitem{openmp-lecture}
    Universitat Aut\`onoma de Barcelona,
    \textit{Lecture 4: OpenMP}, slides 25--37, 44--53.
  \bibitem{starter-code}
    Supplied sequential implementation,
    \href{../../code/XXNN-Multinivell-main/main.c}{\texttt{main.c}},
    functions \texttt{main()}, \texttt{train\_neural\_net()}, and
    \texttt{test\_nn()}.
\end{thebibliography}

\end{document}
